\documentclass[11pt,a4paper]{article}
\usepackage[utf8]{inputenc}
\usepackage[T1]{fontenc}
\usepackage{lmodern}
\usepackage[french]{babel}
\usepackage{amsmath,amssymb,mathtools}
\usepackage{icomma}
\usepackage{xcolor}
\usepackage[most]{tcolorbox}
\usepackage{enumitem}
\usepackage[a4paper,margin=2.2cm,headheight=26pt,headsep=10pt]{geometry}
\usepackage{fancyhdr}
\usepackage[hidelinks]{hyperref}
\definecolor{primary}{RGB}{222,49,99}
\definecolor{primarydark}{RGB}{150,22,55}
\definecolor{excol}{RGB}{0,135,90}
\tcbset{boxbase/.style={enhanced,breakable,boxrule=0.9pt,arc=2.5pt,
  left=3mm,right=3mm,top=2.3mm,bottom=2.3mm,fonttitle=\bfseries,coltitle=white}}
\newtcolorbox[auto counter]{corrige}[1][]{boxbase,colback=excol!4,colframe=excol,
  title={\textsc{Corrigé}~\thetcbcounter\ifx\relax#1\relax\relax\else\textmd{ --- \itshape #1}\fi}}
\newcommand{\R}{\mathbb{R}}
\newcommand{\ds}{\displaystyle}
\pagestyle{fancy}\fancyhf{}
\fancyhead[L]{\small\textcolor{primarydark}{\textbf{Thème 3} — Factorisation et résolution}}
\fancyhead[R]{\small\textcolor{primarydark}{\textsc{Corrigé — Difficile}}}
\fancyfoot[C]{\small\thepage}
\renewcommand{\headrulewidth}{0.8pt}
\begin{document}

\begin{center}
{\LARGE\bfseries\color{primarydark} Niveau Difficile — Corrigé}\\[2pt]
{\color{primary}\rule{0.45\textwidth}{1.2pt}}
\end{center}
\vspace{1mm}

\begin{corrige}[factorisation complète et fraction]
\textbf{a)} $P(1)=4-8+5-1=0$ : $x=1$ est racine. Horner sur $(4,-8,5,-1)$ avec $a=1$ donne les
coefficients $4,-4,1$ et reste $0$, d'où $P(x)=(x-1)(4x^{2}-4x+1)$.

\smallskip
\textbf{b)} $4x^{2}-4x+1=(2x)^{2}-2\cdot2x\cdot1+1^{2}=(2x-1)^{2}$, donc $P(x)=(x-1)(2x-1)^{2}$.

\smallskip
\textbf{c)} $P(x)=0\Leftrightarrow x-1=0$ ou $2x-1=0\Leftrightarrow x=1$ ou $x=\tfrac12$. Donc $S=\{\tfrac12,1\}$.

\smallskip
\textbf{d)} $2x^{2}-3x+1=(2x-1)(x-1)$, C.E. $x\neq1$ et $x\neq\tfrac12$. Alors
\[ \frac{P(x)}{2x^{2}-3x+1}=\frac{(x-1)(2x-1)^{2}}{(2x-1)(x-1)}=2x-1\qquad(x\neq1,\ x\neq\tfrac12). \]
\end{corrige}

\begin{corrige}[trinôme à partir des racines]
\textbf{a)} Deux racines $1$ et $3$ $\Rightarrow P(x)=a(x-1)(x-3)$.

\smallskip
\textbf{b)} $P(2)=a(2-1)(2-3)=a(1)(-1)=-a=2\Rightarrow a=-2$.

\smallskip
\textbf{c)} $P(x)=-2(x-1)(x-3)=-2(x^{2}-4x+3)=-2x^{2}+8x-6$, donc $b=8$ et $c=-6$.
Enfin $P(0)=c=-6$ (ou directement $-2(-1)(-3)=-6$).
\end{corrige}

\begin{corrige}[discriminant et paramètre]
\textbf{a)} Si $m=3$ : l'équation devient $-2x+2=0$, c'est-à-dire $x=1$. C'est une équation du
\emph{premier} degré : \textbf{une seule} solution (le terme en $x^{2}$ disparaît).

\smallskip
\textbf{b)} Pour $m\neq3$ (vrai trinôme, $A=m-3$, $B=-2$, $C=2$) :
\[ \Delta=B^{2}-4AC=(-2)^{2}-4(m-3)(2)=4-8(m-3)=28-8m. \]

\smallskip
\textbf{c)} « Deux racines réelles distinctes » exige un \emph{vrai} trinôme ($m\neq3$) \emph{et} $\Delta>0$ :
\[ 28-8m>0\Leftrightarrow m<\tfrac{28}{8}=\tfrac72. \]
En combinant avec $m\neq3$ : $\ \boxed{\,m\in\,]-\infty,3[\ \cup\ ]3,\tfrac72[\,}$.
\end{corrige}

\end{document}
